\section{Пайплайн подготовки данных}

\subsection{Индексация данных и временнóй порядок}

Этап \texttt{prepare} сканирует каталог с данными, раскладывает найденные файлы по категориям (взаимодействия, пользователи, видео, эмбеддинги) по настраиваемым регулярным выражениям и пишет \texttt{manifest.yaml} --- список файлов, который читают последующие этапы.

Категоризация выполняется по пути \emph{относительно} каталога данных: абсолютный путь может содержать посторонние совпадения с паттернами.

Существенная деталь --- порядок шардов взаимодействий.
Лексикографическая сортировка путей нарушает хронологию: каталог \texttt{test/week\_26} оказывается раньше \texttt{train/week\_00}.
Поэтому шарды сортируются по номеру недели, извлекаемому из имени файла (\texttt{week\_NN}); глобальный номер события \texttt{event\_order} затем присваивается строкам в этом порядке и служит осью времени для сплитов и признаков.

\subsection{Очистка данных и целевая переменная}

Этап \texttt{preprocess} выполняется на Polars в ленивом потоковом режиме: план вычислений строится над \texttt{scan\_parquet}, а материализация происходит сразу в выходной parquet (\texttt{sink\_parquet}), без загрузки всех событий в память.

Шаги очистки:
\begin{enumerate}
    \item отбрасывание событий с пустыми идентификаторами, отрицательным или аномально большим \texttt{timespent} (порог 7200~с), нулевой длительностью видео;
    \item расчет доли досмотра $\text{watch\_ratio} = \text{timespent} / \text{duration}$ с ограничением сверху (клип на 3.0 --- зацикленные пересмотры);
    \item фильтр минимальной активности: не менее 5 событий на пользователя и на видео (агрегаты по более редким сущностям слишком шумные);
    \item опциональное детерминированное прореживание пользователей по хэшу идентификатора (для быстрых прогонов).
\end{enumerate}

На используемой выборке фильтры не отбросили ни одного события (47\,976\,280 строк до и после): выборка была сформирована уже с ограничениями на активность.
На полном датасете эти же фильтры существенны.

Из неявного фидбека строится градуированная целевая переменная (релевантность): досмотр видео и позитивные действия пользователя (лайк, шеринг, закладка, переход к автору, открытие комментариев) увеличивают релевантность, дизлайк обнуляет её.
Градуированный таргет информативнее бинарного и позволяет ранжирующим метрикам различать степень интереса пользователя к видео.

Статистика после этапа: средний \texttt{watch\_ratio} --- 0.456; средняя релевантность --- 0.307; доля позитивных событий ($\text{label} \ge 1$) --- 27.0\%.

\subsection{Признаки}

Этап \texttt{features} строит три группы признаков.
Все агрегаты вычисляются \textbf{только по train-окну} --- первым 80\% событий по \texttt{event\_order}.
Это исключает утечку будущего: признаки объектов валидации и теста опираются лишь на данные, доступные на момент обучения.
Фидбек текущего события (\texttt{like}, \texttt{timespent} и т.\,д.) в признаки не входит --- это компоненты таргета.

\begin{table}[H]
\centering
\caption{Группы признаков (всего 36 числовых и 5 категориальных)}
\label{tab:features}
\small
\begin{tabularx}{\textwidth}{@{}l X@{}}
\toprule
\textbf{Группа} & \textbf{Признаки} \\
\midrule
Пользовательские & статика (\texttt{age}, \texttt{gender}, \texttt{geo}); число событий, средние \texttt{watch\_ratio} и \texttt{timespent}, доли каждого вида фидбека, число уникальных авторов, средняя длительность просматриваемых видео \\
Айтемные & \texttt{duration}; число просмотров, средние \texttt{watch\_ratio} и \texttt{label}, доли фидбека; агрегаты автора: охват, число роликов, средний \texttt{watch\_ratio}, доли фидбека \\
Парные (user--item) & история пользователя с автором видео (\texttt{ua\_*}: число событий, средние \texttt{watch\_ratio} и \texttt{label}); \texttt{dur\_diff} --- отклонение длительности видео от привычной пользователю; \texttt{emb\_cos} --- косинусная близость контентного эмбеддинга видео и эмбеддинг-профиля пользователя \\
Контекст показа & \texttt{place}, \texttt{platform}, \texttt{agent} --- известны на момент ранжирования, используются как категориальные \\
\bottomrule
\end{tabularx}
\end{table}

\subsection{Эмбеддинг-профили пользователей}

Признак \texttt{emb\_cos} вычисляется на PyTorch (устройство выбирается автоматически: CUDA~$\rightarrow$~MPS~$\rightarrow$~CPU) в два потоковых прохода по данным батчами по 500\,000 строк (PyArrow):

\begin{enumerate}
    \item \textbf{Профили.} Для каждого пользователя усредняются L2-нормированные эмбеддинги видео, на которые он дал позитивную реакцию ($\text{label} \ge 1$) в train-окне; профиль также нормируется.
    Аккумуляция выполняется операцией \texttt{index\_add\_} на тензорах.
    \item \textbf{Косинусы.} Для каждого события вычисляется скалярное произведение профиля пользователя и эмбеддинга видео; результат дописывается в parquet инкрементально через \texttt{ParquetWriter}.
\end{enumerate}

Эмбеддинги загружаются из \texttt{npz}-файла с фильтрацией по видео, встречающимся в выборке, и усечением до первых 32 компонент.
На используемой выборке признак получен для 100\% событий; значения лежат в диапазоне $[-0.263;\,0.945]$ при среднем 0.556.

\subsection{Сборка датасетов}

Этап \texttt{build\_dataset} выполняет темпоральное разбиение по \texttt{event\_order} в пропорции 80/10/10 и присоединяет признаки к событиям (таблица~\ref{tab:splits}).

\begin{table}[H]
\centering
\caption{Темпоральное разбиение данных}
\label{tab:splits}
\small
\begin{tabularx}{\textwidth}{@{}l X r@{}}
\toprule
\textbf{Сплит} & \textbf{Диапазон времени} & \textbf{Событий} \\
\midrule
train & первые 80\% событий & 38\,381\,023 \\
validation & следующие 10\% & 4\,797\,628 \\
test & последние 10\% & 4\,797\,629 \\
\bottomrule
\end{tabularx}
\end{table}

Каждый сплит сортируется по паре (\texttt{user\_id}, \texttt{event\_order}): CatBoost требует, чтобы объекты одной группы шли в данных подряд.
Числовые признаки приводятся к \texttt{Float32} (пропуски --- к \texttt{NaN}, которые CatBoost обрабатывает нативно), категориальные --- к строкам с заполнением \texttt{"unknown"}.
Холодные пользователи и видео, впервые появившиеся после train-окна, не отбрасываются --- их агрегаты остаются пропусками.

Списки признаков, размеры сплитов и границы разбиения фиксируются в \texttt{dataset\_meta.yaml}; обучение читает этот файл, что исключает рассинхронизацию между сборкой данных и моделью.
